\documentclass[11pt]{amsart}

% ============================================================
%  Font size + line spacing
% ============================================================
\usepackage{setspace}
\setstretch{1.15}
\setlength{\parskip}{0.4em}
\setlength{\parindent}{1.2em}

\usepackage{amsmath, amssymb, amsthm}
\usepackage[margin=1in,bottom=1cm,footskip=0.7cm]{geometry}
\usepackage{float}
\usepackage{listings}
\usepackage{xcolor}
\usepackage{booktabs}
\usepackage{array}
\usepackage{tikz}
\usetikzlibrary{arrows.meta, positioning}
\usepackage{longtable}
\usepackage{arydshln}
\usepackage{textcase}

% ============================================================
%  Colored headings
% ============================================================
\usepackage{titlesec}
\usepackage{titletoc}

\definecolor{seccolor}{RGB}{0, 51, 102}
\definecolor{subseccolor}{RGB}{0, 80, 130}

\titleformat{\section}
  {\normalfont\Large\bfseries\color{seccolor}}
  {\thesection.}{0.5em}{}
\titleformat{\subsection}
  {\normalfont\large\bfseries\color{subseccolor}}
  {\thesubsection}{0.5em}{}
\titleformat{\subsubsection}
  {\normalfont\normalsize\bfseries\color{subseccolor}}
  {\thesubsubsection}{0.5em}{}

% ============================================================
%  TOC
% ============================================================
\contentsmargin{1.5em}

\titlecontents{section}
  [0pt]
  {\addvspace{0.35em}\small\bfseries\color{seccolor}\raggedright}
  {\contentslabel{2em}}
  {\hspace*{-2em}}
  {\titlerule*[1pc]{.}\contentspage}

\titlecontents{subsection}
  [2em]
  {\small\color{subseccolor}\raggedright}
  {\contentslabel{2.5em}}
  {\hspace*{-2.5em}}
  {\titlerule*[1pc]{.}\contentspage}

\titlecontents{subsubsection}
  [4.5em]
  {\small\color{subseccolor}\raggedright}
  {\contentslabel{3em}}
  {\hspace*{-3em}}
  {\titlerule*[1pc]{.}\contentspage}

% ============================================================
%  Fonts: fall back gracefully if JuliaMono is absent
% ============================================================
\usepackage{iftex}
\ifPDFTeX
  \usepackage[T1]{fontenc}
  \usepackage[utf8]{inputenc}
\else
  \usepackage[no-math]{fontspec}
  \IfFileExists{./JuliaMono/JuliaMono-Regular.ttf}{%
    \setmonofont{JuliaMono}[
      Path = ./JuliaMono/,
      Extension = .ttf,
      UprightFont = *-Regular,
      BoldFont = *-Bold,
      ItalicFont = *-Regular,
      ItalicFeatures = {FakeSlant = 0.2},
      BoldItalicFont = *-Bold,
      BoldItalicFeatures = {FakeSlant = 0.2},
      Scale = 0.85
    ]}{}
\fi

% ============================================================
%  Colored hyperlinks
% ============================================================
\usepackage[
  colorlinks=true,
  linkcolor=seccolor,
  citecolor=teal,
  urlcolor=magenta!80!black,
  filecolor=orange!80!black,
  linktoc=all
]{hyperref}

% ============================================================
%  Custom page numbering
% ============================================================
\usepackage{fancyhdr}
\pagestyle{fancy}
\fancyhf{}
\fancyfoot[C]{\small\color{seccolor}\thepage}
\renewcommand{\headrulewidth}{0pt}
\renewcommand{\footrulewidth}{0pt}

\fancypagestyle{plain}{%
  \fancyhf{}%
  \fancyfoot[C]{\small\color{seccolor}\thepage}%
  \renewcommand{\headrulewidth}{0pt}%
  \renewcommand{\footrulewidth}{0pt}%
}

% ============================================================
%  Lean language for listings
% ============================================================
\lstdefinelanguage{Lean}{
  keywords={import, theorem, have, exact, by, fun, Filter, Tendsto, atTop,
            Nat, Real, nhds, exists, forall, sorry, let, in, if, then, else},
  sensitive=true,
  comment=[l]{--},
  morecomment=[s]{/-}{-/},
  morestring=[b]"
}

% ============================================================
%  Theorem environments
% ============================================================
\theoremstyle{plain}
\newtheorem{axiom}{Axiom}
\newtheorem{definition}{Definition}
\newtheorem{lemma}{Lemma}
\newtheorem{theorem}{Theorem}
\newtheorem{corollary}{Corollary}
\newtheorem{proposition}{Proposition}
\theoremstyle{remark}
\newtheorem{remark}{Remark}

% ============================================================
%  Proof: label bold + underlined, body on new line
% ============================================================
\makeatletter
\renewenvironment{proof}[1][\proofname]{%
  \par\pushQED{\qed}%
  \normalfont\topsep6\p@\@plus6\p@\relax
  \trivlist
  \item[\hskip\labelsep\bfseries\underline{#1}.]%
    \mbox{}\hfill\break
  \ignorespaces
}{%
  \popQED\endtrivlist\@endpefalse
}
\makeatother

% ============================================================
%  Column types for fitchproof
% ============================================================
\newcolumntype{S}{>{\centering\arraybackslash}m{1.0cm}}
\newcolumntype{L}{>{\raggedright\arraybackslash}m{9.0cm}}
\newcolumntype{J}{>{\raggedright\arraybackslash}m{4.0cm}}

% ============================================================
%  Fitch-style proof environment (page-break aware, dashed rules)
%  Single header used for both the first and the continuation heads.
% ============================================================
\newcounter{proofline}

\newenvironment{fitchproof}{%
  \setcounter{proofline}{0}%
  \renewcommand{\arraystretch}{1.2}%
  \setlength{\tabcolsep}{4pt}%
  \setlength{\arrayrulewidth}{0.05pt}%
  \begin{longtable}{|S|L|J|}%
    \hdashline
    \textbf{Step} & \centering{\textbf{Formula}} & \textbf{Justification} \\
    \hdashline
    \endfirsthead
    \hdashline
    \textbf{Step} & \centering{\textbf{Formula}} & \textbf{Justification} \\
    \hdashline
    \endhead
    \hdashline
    \multicolumn{3}{r}{\textit{Continued on next page}} \\
    \endfoot
    \hdashline
    \endlastfoot
  }{%
  \end{longtable}%
}

\newcommand{\pline}[2]{%
  \stepcounter{proofline}%
  \arabic{proofline}. & #1 & \small\textit{#2} \\ \hdashline
}

% ============================================================
%  Listings styles
% ============================================================
\lstdefinestyle{python}{
  language=Python,
  basicstyle=\ttfamily\scriptsize,
  keywordstyle=\color{blue},
  commentstyle=\color{gray},
  stringstyle=\color{red!70!black},
  showstringspaces=false,
  frame=single,
  breaklines=true
}

\lstdefinestyle{lean}{
  language=Lean,
  basicstyle=\ttfamily\scriptsize,
  keywordstyle=\color{blue},
  commentstyle=\color{gray},
  showstringspaces=false,
  frame=single,
  breaklines=true,
  extendedchars=true
}

\raggedbottom

% ============================================================
%  Title, author, address, email, subject class, keywords
% ============================================================
\title[\NoCaseChange{The Diaconis--Mosteller law of truly large numbers}]%
      {\NoCaseChange{The Diaconis--Mosteller law of truly large numbers\\
                    with applications to equidistant letter sequences}}

\author{%
	\NoCaseChange{Brahimi, Mahdi. Tahar. \\
		Department of Mathematics \\
		University of Mohamed Boudiaf, Msila, Algeria \\
		\texttt{\href{mailto:mahdi.brahimi@univ-msila.dz}{mahditahar.brahimi@univ-msila.dz}}
	}
}

\subjclass[2020]{Primary 60C05; Secondary 62E20, 68R15, 60F05}
\keywords{Law of truly large numbers, coincidence, equidistant letter
sequence, Bible code, multiple testing, Poisson approximation,
combinatorial probability, first moment method, second moment method}

\begin{document}

\begin{abstract}
	We formalize the \emph{law of truly large numbers}, a statistical
	observation attributed to Persi Diaconis and Frederick Mosteller \cite{diaconis1989}. The law
	states that with a sufficiently large number of independent trials, any
	event with constant positive probability, however small, is almost certain
	to occur at least once. We present the foundational axioms, definitions,
	lemmas, theorems, propositions, and corollaries that constitute the
	mathematical framework of this law, along with complete proofs. We also
	apply the law to Equidistant Letter Sequences (ELS) in long texts, showing
	that the ``Bible code'' phenomenon is a direct consequence of the law of
	truly large numbers. All combinatorial constants are stated exactly: the
	number of candidate ELS pairs is
	\(T = N^{2}/(L-1) + O(N)\), not the frequently quoted \(2N^{2}\).
\end{abstract}

\maketitle

\begingroup
  \small
  \setlength{\parskip}{0pt}
  \setcounter{tocdepth}{2}
  \tableofcontents
\endgroup
\newpage

% ============================================================
\section{Axioms}
% ============================================================

\begin{axiom}[Independence]
\label{ax:independence}
Events in a sequence of trials are independent if the occurrence of any
event does not affect the probability of any other event.
\end{axiom}

\begin{axiom}[Constant Probability]
\label{ax:constant}
There exists a fixed probability \(p \in (0,1]\) such that in each trial,
the event \(A\) occurs with probability \(p\), independently of all other
trials.
\end{axiom}

% ============================================================
\section{Definitions}
% ============================================================

\begin{definition}[Rare Event]
\label{def:rare}
An event \(A\) is called \emph{rare} if its probability \(p\) in a single
trial is very small, i.e., \(p \ll 1\).
\end{definition}

\begin{definition}[Truly Large Number of Trials]
\label{def:large}
A number of trials \(N\) is called \emph{truly large} if \(N\) is
sufficiently large that the expected number of occurrences, \(Np\), is
large, or more precisely, if \((1-p)^N\) is close to zero.
\end{definition}

\begin{definition}[Law of Truly Large Numbers]
\label{def:law}
The law of truly large numbers states that for any fixed \(p \in (0,1]\),
the probability that \(A\) occurs at least once in \(N\) independent
trials tends to \(1\) as \(N \to \infty\).
\end{definition}

\begin{definition}[Text and Alphabet]
\label{def:text}
Let \(S = s_1 s_2 \cdots s_N\) be a finite string of length \(N\) over an
alphabet \(\mathcal{A}\) of size \(q = |\mathcal{A}|\). We assume the
letters are drawn independently and uniformly from \(\mathcal{A}\) under
the null hypothesis.
\end{definition}

\begin{definition}[Equidistant Letter Sequence]
\label{def:els}
An \emph{Equidistant Letter Sequence} (ELS) for a word
\(w = w_1 w_2 \cdots w_L\) is a pair \((i,d)\) with \(1 \le i \le N\),
\(d \neq 0\), and \(1 \le i + (L-1)d \le N\), such that
\begin{equation}
\label{eq:els-def}
s_{i + k d} = w_{k+1}, \qquad k = 0, 1, \dots, L-1.
\end{equation}
The integer \(d\) is called the \emph{skip}. Positive \(d\) corresponds to
forward reading; negative \(d\) to backward reading.
\end{definition}

\begin{definition}[Occurrence Count]
\label{def:count}
Let \(X_w\) denote the total number of ELS occurrences of the word \(w\)
in \(S\):
\begin{equation}
\label{eq:count}
X_w = \#\{(i,d) : (i,d) \text{ is an ELS for } w\}.
\end{equation}
Equivalently \(X_w = \sum_{u} X_u\) where \(X_u\) is the indicator that
candidate \(u=(i,d)\) is a hit.
\end{definition}

\begin{definition}[Position Set]
\label{def:positions}
For a candidate \(u = (i,d)\), the \emph{position set} is
\begin{equation}
\label{eq:positions}
P_u := \{\,i,\; i+d,\; i+2d,\; \dots,\; i+(L-1)d\,\} \subseteq \{1,\dots,N\}.
\end{equation}
\end{definition}

\begin{definition}[Aperiodic Word]
\label{def:aperiodic}
A word \(w = w_1 \cdots w_L\) is \emph{aperiodic} if it has no nontrivial
period: there is no \(j \in \{1,\dots,L-1\}\) such that
\(w_k = w_{k+j}\) for all \(k \in \{1,\dots,L-j\}\). Otherwise \(w\) is
\emph{periodic}.
\end{definition}

\begin{definition}[Candidate-Pair Set and Neighbourhood]
\label{def:candidates}
Let \(I\) denote the set of all admissible candidates \((i,d)\),
\(T=|I|\). For \(u\in I\) let
\begin{equation}
\label{eq:Bu-def}
B_u := \{\,v=(i',d')\in I : v\neq u,\ P_u\cap P_v \neq \varnothing\,\}
\end{equation}
be the \emph{overlap neighbourhood} of \(u\).
\end{definition}

\begin{definition}[Multiple Testing]
\label{def:mult}
If one searches for \(M\) distinct words or spelling variants, all of
length \(L\), the total expected number of hits is
\begin{equation}
\label{eq:mult-test}
\mathbb{E}[X_{\text{total}}] \approx M \cdot \frac{N^{2}}{(L-1)\,q^{L}}.
\end{equation}
\end{definition}

% ============================================================
\section{Statements}
% ============================================================

\subsection{Lemmas}

\begin{lemma}[Probability of No Occurrence]
\label{lem:noocc}
Let \(A\) be an event with probability \(p\) in a single trial. The
probability that \(A\) does not occur in \(N\) independent trials is
\begin{equation}
\label{eq:no-occ}
\Pr(A^c \text{ in all } N \text{ trials}) = (1-p)^N.
\end{equation}
\end{lemma}

\begin{lemma}[Number of Candidate Pairs]
\label{lem:candidates}
Put \(D := \lfloor (N-1)/(L-1) \rfloor\). Let \(T\) be the number of pairs
\((i,d)\) with \(1 \le i \le N\), \(d \neq 0\), and
\(1 \le i + (L-1)d \le N\). Then, with \(L\) fixed and \(N \to \infty\),
\begin{equation}
\label{eq:candidates-exact}
T \;=\; 2\!\left[ DN - \frac{(L-1)\,D(D+1)}{2} \right]
   \;=\; \frac{N^{2}}{L-1} + O(N).
\end{equation}
Equivalently, \(T = 2\,T_{D}(N,L)\) where
\(T_{D}(N,L) = \sum_{d=1}^{D}\bigl(N-(L-1)d\bigr)\) is the triangular sum.
\end{lemma}

\begin{lemma}[Size of the Overlap Neighbourhood]
\label{lem:Bu}
For every \(u\in I\),
\begin{equation}
\label{eq:Bu-bound}
|B_u| \;\le\; \frac{2\,L^{2}\,N}{L-1}.
\end{equation}
Consequently
\begin{equation}
\label{eq:sum-Bu}
\sum_{u\in I} |B_u| \;\le\; \frac{2\,L^{2}\,N^{3}}{(L-1)^{2}}.
\end{equation}
\end{lemma}

\begin{lemma}[Second Moment of the ELS Count]
\label{lem:variance}
Under the null hypothesis of Definition~\ref{def:text},
\begin{equation}
\label{eq:second-moment}
\mathbb{E}[X_w(X_w-1)]
  = \sum_{\substack{u,v\in I\\ u\neq v}} \Pr(X_u = X_v = 1).
\end{equation}
Partition the ordered pairs \((u,v)\), \(u\neq v\), into
\emph{disjoint} pairs (\(P_u\cap P_v=\varnothing\)) and
\emph{overlapping} pairs (\(P_u\cap P_v\neq\varnothing\)).
Then
\begin{equation}
\label{eq:aperiodic-cov}
\mathbb{E}[X_w(X_w-1)]
   \;=\; \lambda_w^{2}\bigl(1 + o(1)\bigr) + O\!\left(\frac{N^{3}}{q^{2L}}\right),
\qquad
\lambda_w = \frac{N^{2}}{(L-1)\,q^{L}},
\end{equation}
and hence
\begin{equation}
\label{eq:variance}
\operatorname{Var}(X_w)
   \;=\; \mathbb{E}[X_w] + O\!\left(\frac{N^{3}}{q^{2L}}\right).
\end{equation}
\end{lemma}

\subsection{Theorems}

\begin{theorem}[Law of Truly Large Numbers]
\label{thm:law}
Let \(A\) be an event with probability \(p \in (0,1]\). Then
\begin{equation}
\label{eq:law-limit}
\lim_{N \to \infty} \Pr(A \text{ occurs at least once in } N \text{ trials}) = 1.
\end{equation}
\end{theorem}

\begin{theorem}[Poisson Approximation, Independent Case]
\label{thm:poisson}
For rare events with large \(N\) and small \(p\) such that \(Np = \lambda\)
is constant, the number of occurrences \(X\) converges in distribution to
a Poisson random variable with parameter \(\lambda\):
\begin{equation}
\label{eq:poisson-pmf}
\Pr(X = k) \approx \frac{e^{-\lambda} \lambda^k}{k!}.
\end{equation}
In particular, \(\Pr(X \ge 1) \approx 1 - e^{-\lambda}\).
\end{theorem}

\begin{theorem}[Expected ELS Hits for Uniform Text]
\label{thm:els}
Under the null hypothesis of independent uniform letters, the expected
number of ELS occurrences of a word \(w\) of length \(L\) satisfies
\begin{equation}
\label{eq:els-moment}
\mathbb{E}[X_w]
  = \frac{1}{q^{L}}\sum_{d \neq 0}\max\!\bigl(0,\,N-(L-1)|d|\bigr)
  = \frac{N^{2}}{(L-1)\,q^{L}} + O\!\left(\frac{N}{q^{L}}\right).
\end{equation}
In particular, for \(L \ll N\),
\begin{equation}
\label{eq:els-moment-asymp}
\mathbb{E}[X_w] \approx \frac{N^{2}}{(L-1)\,q^{L}}.
\end{equation}
\end{theorem}

\begin{theorem}[Poisson Approximation for Aperiodic Words]
\label{thm:poisson-els}
Let \(w\) be aperiodic of length \(L\), and suppose \(N \to \infty\) with
\(L\) and \(q\) fixed in the \emph{sparse regime} \(N/q^{L} \to 0\). Then
\begin{equation}
\label{eq:poisson-els}
d_{TV}\!\left(X_w,\ \mathrm{Pois}(\lambda_w)\right) \;\xrightarrow[N\to\infty]{}\; 0,
\qquad
\lambda_w = \frac{N^{2}}{(L-1)\,q^{L}}.
\end{equation}
\end{theorem}

\subsection{Propositions}

\begin{proposition}[Expected Number of Occurrences]
\label{prop:expected}
The expected number of occurrences of \(A\) in \(N\) trials is \(Np\). If
\(Np \ge 1\), the event is expected to occur at least once.
\end{proposition}

\begin{proposition}[First Moment Method]
\label{prop:first-moment}
Let \(X \ge 0\) be an integer-valued random variable. If
\(\mathbb{E}[X] \to \infty\), then \(\Pr(X \ge 1) \to 1\).
\end{proposition}

\begin{proposition}[Birthday Problem as a Special Case]
\label{prop:birthday}
In a group of \(N\) people, the probability that at least two share a
birthday is
\begin{equation}
\label{eq:birthday}
1 - \frac{365!}{(365-N)! \, 365^N}.
\end{equation}
For \(N=23\), this probability exceeds \(0.5\).
\end{proposition}

\begin{proposition}[Condition for Expected Hits]
\label{prop:condition}
The expected number of ELS hits exceeds \(1\) when
\begin{equation}
\label{eq:threshold}
N \gtrsim q^{L/2}\sqrt{L-1}.
\end{equation}
For the Hebrew Torah with \(N \approx 304{,}805\) and \(q = 22\), this
gives the following expected counts:
\begin{longtable}{c|c|c}
	\caption{Expected ELS counts for the Hebrew Torah with
		\(N \approx 304{,}805\) and \(q = 22\).}
	\label{tab:torah}\\
	\toprule
	\(L\)
	& \(\mathbb{E}[X_w] \approx \dfrac{N^{2}}{(L-1)\,22^{L}}\)
	& \(\Pr(X_w \ge 1) \approx 1 - e^{-\mathbb{E}[X_w]}\) \\
	\midrule
	\endfirsthead
	
	\toprule
	\(L\)
	& \(\mathbb{E}[X_w] \approx \dfrac{N^{2}}{(L-1)\,22^{L}}\)
	& \(\Pr(X_w \ge 1) \approx 1 - e^{-\mathbb{E}[X_w]}\) \\
	\midrule
	\endhead
	
	\midrule
	\multicolumn{3}{r}{\textit{Continued on next page}} \\
	\endfoot
	
	\bottomrule
	\endlastfoot
	
	5  & \(4.5 \times 10^{3}\)   & \(> 1 - 10^{-1900}\) \\
	6  & \(1.6 \times 10^{2}\)   & \(> 1 - 10^{-69}\) \\
	7  & \(6.2\)                 & \(0.998\) \\
	8  & \(0.24\)                & \(0.213\) \\
	9  & \(9.6 \times 10^{-3}\)  & \(0.0096\) \\
	10 & \(3.9 \times 10^{-4}\)  & \(3.9 \times 10^{-4}\) \\
\end{longtable}
\end{proposition}

\begin{proposition}[Look-Elsewhere Effect]
\label{prop:lookelsewhere}
A nominal \(p\)-value \(p_0\) for a single ELS hit must be corrected by the
effective number of tests \(T_{\text{eff}}\) to obtain a valid significance
level:
\begin{equation}
\label{eq:lookelsewhere}
p_{\text{corrected}} \approx 1 - (1 - p_0)^{T_{\text{eff}}} \approx p_0 \cdot T_{\text{eff}}.
\end{equation}
For ELS searches, \(T_{\text{eff}}\) is on the order of \(M\,N^{2}/(L-1)\),
which is typically enormous.
\end{proposition}

\begin{proposition}[Bible Code Hits Are Expected]
\label{prop:bible}
For the Hebrew Torah (\(N \approx 304{,}805\), \(q = 22\)), any 5--7
letter word is expected to appear many times by chance. Names such as
``Clinton,'' ``Rabin,'' ``Sharon,'' and ``Arafat'' are therefore expected
to be found in ELS searches without any hidden design.
\end{proposition}

\begin{proposition}[McKay's Rebuttal]
\label{prop:mckay}
Brendan McKay and collaborators demonstrated that similar ELS patterns can
be found in \emph{Moby Dick} and \emph{War and Peace}, confirming that the
phenomenon is generic to long texts and not unique to the Bible.
\end{proposition}

\subsection{Corollaries}

\begin{corollary}[Quantitative Bound]
\label{cor:bound}
For a given significance level \(\alpha \in (0,1)\), the event \(A\) occurs
with probability at least \(1-\alpha\) whenever the (integer) number of
trials \(N\) satisfies
\begin{equation}
\label{eq:min-trials}
N \ge N_{\min}(\alpha,p) := \left\lceil \frac{\ln \alpha}{\ln(1-p)} \right\rceil .
\end{equation}
\end{corollary}

\begin{corollary}[Second Moment Method]
\label{cor:second-moment}
Let \(X \ge 0\) be integer-valued with \(\mathbb{E}[X] = \mu > 0\) and
\(\operatorname{Var}(X) = \sigma^{2}\). Then
\begin{equation}
\label{eq:second-moment-cor}
\Pr(X \ge 1) \;\ge\; \frac{\mu^{2}}{\mu^{2} + \sigma^{2}}.
\end{equation}
In particular, for the ELS count \(X_w\) with aperiodic \(w\) and
\(\lambda_w \to \infty\) while \(\sigma^{2} = \lambda_w\,(1+o(1))\),
\begin{equation}
\label{eq:els-second}
\Pr(X_w \ge 1) \;\ge\; \frac{\lambda_w}{\lambda_w+1} \to 1.
\end{equation}
\end{corollary}

\begin{corollary}[Coincidences Are Expected]
\label{cor:coincidence}
In a population of size \(M\), if each pair of individuals has a small
probability \(q\) of sharing some rare characteristic, the expected number
of coincident pairs is \(\binom{M}{2} q\). For large \(M\), this expectation
can be substantial even if \(q\) is tiny.
\end{corollary}

\begin{corollary}[Non-Uniform Letter Frequencies]
\label{cor:nonuniform}
If letter \(a \in \mathcal{A}\) has frequency \(p_a\), then
\begin{equation}
\label{eq:nonuniform}
\mathbb{E}[X_w] \approx \frac{N^{2}}{L-1}\prod_{k=1}^L p_{w_k}.
\end{equation}
\end{corollary}

\begin{corollary}[Invalidity of Naive Odds]
\label{cor:invalid}
Reported odds such as \(10^5:1\) or \(10^6:1\) for ELS hits are almost
always uncorrected single-test \(p\)-values. After correcting for the
number of trials, the significance vanishes.
\end{corollary}

\begin{corollary}[No Predictive Power]
\label{cor:nopredict}
ELS searches cannot reliably predict future events, because the apparent
hits are artifacts of the law of truly large numbers and post-hoc
selection.
\end{corollary}

\begin{remark}[Why the disjoint/overlapping split is correct]
\label{rem:correct-split}
For aperiodic \(w\), same-skip shifts \(v=(i+jd,d)\), \(1\le j\le L-1\),
force \(w\) to be \(j\)-periodic; since \(w\) is aperiodic, the joint
probability of \((X_u,X_v)=(1,1)\) vanishes identically. For different
skips \(d'\neq d\), overlapping pairs need not have vanishing
probability: they can only contribute through position-set intersections
of size \(m\ge 1\), and each such pair has probability at most
\(q^{-(2L-1)}\). Since the number of overlapping pairs is bounded by
\(\sum_u |B_u| \le 2L^{2}N^{3}/(L-1)^{2}\) (Lemma~\ref{lem:Bu}), the total
overlapping contribution is
\(O\!\bigl(N^{3}q^{-(2L-1)}/(L-1)\bigr) = O(N^{3}/q^{2L})\). This is the
source of the error term in Lemma~\ref{lem:variance}. The leading term
\(\lambda_w^{2}\) comes from the disjoint pairs, not from the diagonal.
\end{remark}

\begin{remark}[Chen--Stein bound for aperiodic words]
\label{rem:chenstein}
Let \(p = q^{-L}\) and \(b_1,b_2\) be as in the Chen--Stein total-variation
bound for locally dependent indicators:
\begin{equation}
\label{eq:CSbound}
d_{TV}\!\left(X_w, \mathrm{Pois}(\lambda_w)\right)
\;\le\; \frac{1-e^{-\lambda_w}}{\lambda_w}\,(b_1 + b_2),
\qquad
b_1 = \sum_{u}\sum_{v\in B_u} p^{2},
\qquad
b_2 = \sum_{u}\sum_{v\in B_u} p_{uv}.
\end{equation}
For aperiodic \(w\), the same-skip contributions to \(b_2\) vanish
(Remark~\ref{rem:correct-split}), and the different-skip contributions
are bounded by \(p_{uv} \le q^{-(2L-1)}\). Hence
\begin{equation}
\label{eq:CS-normalized}
\frac{b_1}{\lambda_w} = O\!\left(\frac{N}{q^{L}}\right),
\qquad
\frac{b_2}{\lambda_w} = O\!\left(\frac{N}{q^{L}}\right).
\end{equation}
Therefore \(d_{TV}\to 0\) whenever \(N/q^{L}\to 0\), which is the sparse
regime of Theorem~\ref{thm:poisson-els}. For highly periodic \(w\), the
naive Chen--Stein bound degenerates (since \(b_2\asymp\lambda_w\));
this is a genuine feature of the ELS dependence structure and requires
either the exponential-dependency-graph inequalities of Janson
\cite{janson2004} or a direct second-moment argument. We leave that
refinement to future work.
\end{remark}

% ============================================================
\section{Proofs}
% ============================================================

\begin{proof}[Proof of Lemma \ref{lem:noocc}]
  \begin{fitchproof}
    \pline{\(N\) independent trials with \(\Pr(A_i^c) = 1-p\)}{premise}
    \pline{\(\Pr\!\left(\bigcap_{i=1}^N A_i^c\right) = \prod_{i=1}^N \Pr(A_i^c)\)}{independence (Axiom \ref{ax:independence})}
    \pline{\(= \prod_{i=1}^N (1-p)\)}{substitution}
    \pline{\(= (1-p)^N\)}{algebra, gives\, Eq. ~\eqref{eq:no-occ}}
  \end{fitchproof}
\end{proof}

\begin{proof}[Proof of Lemma \ref{lem:candidates}]
  \begin{fitchproof}
    \pline{Fix \(d>0\). Valid starts: \(i \le N-(L-1)d\).}{counting}
    \pline{Let \(D = \lfloor (N-1)/(L-1)\rfloor\). Then}
    {short-hands}
    \pline{\(T_+ = \sum_{d=1}^{D} \bigl(N-(L-1)d\bigr)
      = DN - (L-1)\tfrac{D(D+1)}{2}\).}{arithmetic series}
    \pline{Write \(D = \tfrac{N-1}{L-1} + O(1)\); expanding,}
    {substitution}
    \pline{\(T_+ = \tfrac{N^{2}}{2(L-1)} + O(N)\).}{simplify}
    \pline{By symmetry \(d<0\) contributes \(T_- = T_+\).}{symmetry}
    \pline{\(T = 2T_+ = \tfrac{N^{2}}{L-1} + O(N)\).}{sum, gives\, Eq. ~\eqref{eq:candidates-exact}}
  \end{fitchproof}
\end{proof}

\begin{proof}[Proof of Lemma \ref{lem:Bu}]
  \begin{fitchproof}
    \pline{Fix \(u=(i,d)\) and \(v=(i',d')\) with \(v\in B_u\),
      so there is \(p\in P_u\cap P_v\).}{premise}
    \pline{Write \(p = i + k d = i' + k' d'\) for some
      \(k,k'\in\{0,\dots,L-1\}\).}{definition of \(P_u\)}
    \pline{For fixed \(k,d',p\), the value \(i'=p-k'd'\) is determined
      for each \(k'\in\{0,\dots,L-1\}\).}{linear algebra}
    \pline{Hence for fixed \(u\) and fixed \(d'\), at most \(L^{2}\)
      values of \(i'\) yield \(v\in B_u\).}{counting}
    \pline{There are at most \(\lceil 2N/(L-1)\rceil\) nonzero values
      of \(d'\) with \(|d'|\le N/(L-1)\).}{range of \(d'\)}
    \pline{Therefore \(|B_u| \le L^{2}\cdot 2N/(L-1)
      = 2L^{2}N/(L-1)\).}{product, gives\, Eq. ~\eqref{eq:Bu-bound}}
    \pline{Summing over \(u\in I\) and using \(T \le N^{2}/(L-1)+O(N)\):}
    {Lemma \ref{lem:candidates}}
    \pline{\(\sum_u |B_u| \le 2L^{2}N^{3}/(L-1)^{2} + O(L^{2}N^{2}/(L-1))\).}{summation}
  \end{fitchproof}
\end{proof}

\begin{proof}[Proof of Lemma \ref{lem:variance}]
  \begin{fitchproof}
    \pline{\(X_w(X_w-1) = \sum_{u\neq v} X_u X_v\).}{indicator identity}
    \pline{\(\mathbb{E}[X_w(X_w-1)] = \sum_{u\neq v} \Pr(X_u=X_v=1)\).}{linearity of expectation}
    \pline{Partition the ordered pairs: disjoint \(D\) and overlapping \(O\).}{combinatorial split}
    \pline{For \(u\neq v\) disjoint, \(\Pr(X_u=X_v=1) = q^{-2L}\),
      independent of \(w\).}{independence of position sets}
    \pline{Number of disjoint ordered pairs:
      \(|D| = T(T-1) - |O|\).}{counting}
    \pline{By Lemma~\ref{lem:Bu}, \(|O| \le \sum_u |B_u|
      \le 2L^{2}N^{3}/(L-1)^{2}\).}{bound on \(O\)}
    \pline{Hence \(|D| = T^{2}\bigl(1 + O(L^{2}/(N))\bigr)\).}{asymptotic expansion}
    \pline{Disjoint contribution:
      \(|D|\,q^{-2L} = \lambda_w^{2}\bigl(1+O(1/N)\bigr)\).}{first term}
    \pline{Overlapping contribution: for each \((u,v)\in O\),
      \(\Pr(X_u=X_v=1) \le q^{-(2L-1)}\).}{worst-case bound}
    \pline{Hence overlapping contribution is
      \(O(|O|\,q^{-(2L-1)}) = O(N^{3}/q^{2L})\).}{substitution, gives\, Eq. ~\eqref{eq:aperiodic-cov}}
    \pline{Combine: \(\mathbb{E}[X_w(X_w-1)]
      = \lambda_w^{2}(1+o(1)) + O(N^{3}/q^{2L})\).}{sum}
    \pline{Variance identity:
      \(\operatorname{Var}(X_w) = \mathbb{E}[X_w(X_w-1)] + \mathbb{E}[X_w]
      - \mathbb{E}[X_w]^{2}\).}{definition of variance}
    \pline{\(= \lambda_w^{2}(1+o(1)) + O(N^{3}/q^{2L})
      + \lambda_w - \lambda_w^{2}(1+o(1))\).}{substitute}
    \pline{\(= \lambda_w + O(N^{3}/q^{2L})\).}{cancel leading terms, gives\, Eq. ~\eqref{eq:variance}}
  \end{fitchproof}
\end{proof}

\begin{proof}[Proof of Theorem \ref{thm:law}]
  \begin{fitchproof}
    \pline{\(\Pr(A \text{ at least once}) = 1 - \Pr(A^c \text{ in all } N)\)}{complement}
    \pline{\(= 1 - (1-p)^N\)}{by\, Eq. ~\eqref{eq:no-occ}}
    \pline{Since \(0 < p \le 1\), \(0 \le 1-p < 1\).}{arithmetic}
    \pline{\((1-p)^N \to 0\).}{geometric decay}
    \pline{\(\therefore\; 1 - (1-p)^N \to 1\).}{limit, gives\, Eq. ~\eqref{eq:law-limit}}
  \end{fitchproof}
\end{proof}

\begin{proof}[Proof of Theorem \ref{thm:poisson}]
  \begin{fitchproof}
    \pline{Binomial pmf: \(\displaystyle\Pr(X=k) = \binom{N}{k} p^{k} (1-p)^{N-k}\).}{definition}
    \pline{Take \(N \to \infty\), \(p \to 0\), \(Np \to \lambda\).}{law of rare events}
    \pline{\(\binom{N}{k} p^{k} (1-p)^{N-k} \to \dfrac{e^{-\lambda}\lambda^{k}}{k!}\).}{Poisson limit theorem}
  \end{fitchproof}
\end{proof}

\begin{proof}[Proof of Theorem \ref{thm:els}]
  \begin{fitchproof}
    \pline{Per candidate pair \((i,d)\): \(\Pr(\text{match}) = q^{-L}\).}{uniform alphabet}
    \pline{\(\mathbb{E}[X_w] = \sum_{(i,d)} \Pr(\text{match}_{(i,d)})\).}{linearity of expectation}
    \pline{\(= \sum_{d \neq 0} \max(0,N-(L-1)|d|)\,q^{-L}\).}{by step 1}
    \pline{\(= T q^{-L}\).}{definition of \(T\)}
    \pline{\(= \dfrac{N^{2}}{(L-1)q^{L}} + O\!\left(\dfrac{N}{q^{L}}\right)\).}{by\, Eq. ~\eqref{eq:candidates-exact}, gives\, Eq. ~\eqref{eq:els-moment}}
  \end{fitchproof}
\end{proof}

\begin{proof}[Proof of Theorem \ref{thm:poisson-els}]
  \begin{fitchproof}
    \pline{Use Chen--Stein bound\, Eq. ~\eqref{eq:CSbound}.}{Remark~\ref{rem:chenstein}}
    \pline{For aperiodic \(w\), \(b_1/\lambda_w = O(N/q^{L})\)
      and \(b_2/\lambda_w = O(N/q^{L})\).}{Remark~\ref{rem:chenstein}, \, Eq. ~\eqref{eq:CS-normalized}}
    \pline{\((1-e^{-\lambda_w})/\lambda_w \le 1\).}{elementary bound}
    \pline{\(d_{TV}(X_w,\mathrm{Pois}(\lambda_w)) = O(N/q^{L})\).}{substitute}
    \pline{\(N/q^{L} \to 0\) by hypothesis, so \(d_{TV} \to 0\).}{limit, gives\, Eq. ~\eqref{eq:poisson-els}}
  \end{fitchproof}
\end{proof}

\begin{proof}[Proof of Proposition \ref{prop:expected}]
  \begin{fitchproof}
    \pline{\(X \sim \mathrm{Binomial}(N, p)\).}{definition}
    \pline{\(\mathbb{E}[X] = Np\).}{standard}
    \pline{If \(Np \ge 1\), \(\mathbb{E}[X] \ge 1\).}{immediate}
  \end{fitchproof}
\end{proof}

\begin{proof}[Proof of Proposition \ref{prop:first-moment}]
  \begin{fitchproof}
    \pline{Markov's inequality: \(\Pr(X \ge 1) \ge \Pr(X > 0)\).}{elementary}
    \pline{\(\Pr(X=0) = 1 - \Pr(X \ge 1) \le 1 - \mu/(1+\mu)\).}{elementary}
    \pline{Alternatively: \(\Pr(X \ge 1) = 1 - \Pr(X=0)\).}{complement}
    \pline{\(\Pr(X=0) = \Pr(X \le 0) \le \mathbb{E}[e^{-tX}]\) for \(t>0\).}{Chernoff}
    \pline{For integer \(X \ge 0\): \(\Pr(X=0) \le 1/(1+\mathbb{E}[X])\).}{conditional expectation}
    \pline{Hence \(\Pr(X \ge 1) \ge \mathbb{E}[X]/(1+\mathbb{E}[X]) \to 1\).}{limit, gives Proposition}
  \end{fitchproof}
\end{proof}

\begin{proof}[Proof of Proposition \ref{prop:birthday}]
  \begin{fitchproof}
    \pline{\(\Pr(\text{all distinct}) = \prod_{k=0}^{N-1} \frac{365-k}{365}\).}{sequential}
    \pline{\(= \dfrac{365!}{(365-N)! \, 365^N}\).}{algebra}
    \pline{\(\Pr(\ge 1 \text{ match}) = 1 - \Pr(\text{all distinct})\).}{complement, gives\, Eq. ~\eqref{eq:birthday}}
    \pline{For \(N=23\), numerical evaluation gives \(> 0.5\).}{numerical}
  \end{fitchproof}
\end{proof}

\begin{proof}[Proof of Proposition \ref{prop:condition}]
  \begin{fitchproof}
    \pline{\(\mathbb{E}[X_w] \approx N^{2}/((L-1)q^{L})\).}{by\, Eq. ~\eqref{eq:els-moment-asymp}}
    \pline{Threshold: \(\mathbb{E}[X_w] > 1 \iff N^{2} > (L-1)q^{L}\).}{algebra}
    \pline{\(\iff N > q^{L/2}\sqrt{L-1}\).}{solve for \(N\), gives\, Eq. ~\eqref{eq:threshold}}
    \pline{At \(N=304{,}805\), \(q=22\): substitute into \(N^{2}/((L-1)22^{L})\).}{numeric, gives\, table. ~\ref{tab:torah}}
  \end{fitchproof}
\end{proof}

\begin{proof}[Proof of Proposition \ref{prop:lookelsewhere}]
  \begin{fitchproof}
    \pline{\(\Pr(\ge 1 \text{ hit among } T_{\text{eff}} \text{ tests}) = 1 - (1-p_0)^{T_{\text{eff}}}\).}{complement}
    \pline{\(\approx 1 - (1 - p_0 T_{\text{eff}})\) for small \(p_0\).}{binomial approx.}
    \pline{\(= p_0 T_{\text{eff}}\).}{algebra, gives\, Eq. ~\eqref{eq:lookelsewhere}}
    \pline{Here \(T_{\text{eff}} \asymp M N^{2}/(L-1)\).}{by\, Eq. ~\eqref{eq:els-moment}}
  \end{fitchproof}
\end{proof}

\begin{proof}[Proof of Proposition \ref{prop:bible}]
  \begin{fitchproof}
    \pline{For Torah \(N \approx 304{,}805\), \(q=22\), \(\mathbb{E}[X_w]\) from\, table ~\ref{tab:torah}.}{numerical}
    \pline{For \(L=5,6,7\): \(\mathbb{E}[X_w] = 4.5\!\times\!10^{3},\ 1.6\!\times\!10^{2},\ 6.2\).}{all \(\gg 1\)}
    \pline{\(\Pr(X_w \ge 1) \to 1\) for any 5--7 letter word.}{Poisson heuristic}
    \pline{Names ``Clinton,'' ``Rabin,'' ``Sharon,'' ``Arafat'' are 5--7 letters.}{direct application}
    \pline{Such hits are expected without hidden design.}{conclusion}
  \end{fitchproof}
\end{proof}

\begin{proof}[Proof of Proposition \ref{prop:mckay}]
  \begin{fitchproof}
    \pline{McKay et al.\ search \emph{Moby Dick} and \emph{War and Peace} for ELS patterns.}{see~\cite{mckay1999}}
    \pline{Similar patterns appear in both non-biblical texts.}{empirical finding}
    \pline{ELS phenomenon is generic to long texts.}{conclusion}
    \pline{It is not unique to the Bible.}{conclusion}
  \end{fitchproof}
\end{proof}

\begin{proof}[Proof of Corollary \ref{cor:bound}]
  \begin{fitchproof}
    \pline{Goal: \(1 - (1-p)^N \ge 1-\alpha\).}{target}
    \pline{\(\iff (1-p)^N \le \alpha\).}{algebra}
    \pline{\(\iff N \ln(1-p) \le \ln \alpha\).}{take \(\ln\)}
    \pline{\(\iff N \ge \ln\alpha/\ln(1-p)\).}{divide by \(\ln(1-p)<0\)}
    \pline{Integer \(N\): ceiling.}{rounding, gives\, Eq. ~\eqref{eq:min-trials}}
  \end{fitchproof}
\end{proof}

\begin{proof}[Proof of Corollary \ref{cor:second-moment}]
  \begin{fitchproof}
    \pline{\(\Pr(X \ge 1) \ge \mu^{2}/(\mu^{2}+\sigma^{2})\).}{Paley--Zygmund inequality applied to \(X>0\)}
    \pline{For ELS \(X_w\): \(\mu = \lambda_w\), \(\sigma^{2} = \lambda_w(1+o(1))\).}{Lemma~\ref{lem:variance}}
    \pline{\(\Pr(X_w \ge 1) \ge \lambda_w/(\lambda_w+1) \to 1\).}{limit, gives\, Eq. ~\eqref{eq:els-second}}
  \end{fitchproof}
\end{proof}

\begin{proof}[Proof of Corollary \ref{cor:coincidence}]
  \begin{fitchproof}
    \pline{Number of pairs \(= \binom{M}{2}\).}{combinatorics}
    \pline{\(\Pr(\text{match per pair}) = q\).}{given}
    \pline{\(\mathbb{E}[\#\text{matches}] = \binom{M}{2}q\).}{linearity of expectation}
  \end{fitchproof}
\end{proof}

\begin{proof}[Proof of Corollary \ref{cor:nonuniform}]
  \begin{fitchproof}
    \pline{\(\Pr(\text{match}) = \prod_{k=1}^L p_{w_k}\).}{independence}
    \pline{\(\mathbb{E}[X_w] = T \prod_{k=1}^L p_{w_k}\).}{linearity}
    \pline{\(= \dfrac{N^{2}}{L-1} \prod_{k=1}^L p_{w_k} + O(N)\).}{by\, Eq. ~\eqref{eq:candidates-exact}, gives\, Eq. ~\eqref{eq:nonuniform}}
  \end{fitchproof}
\end{proof}

\begin{proof}[Proof of Corollary \ref{cor:invalid}]
  \begin{fitchproof}
    \pline{Nominal odds: single-test \(p_0\).}{premise}
    \pline{Corrected: \(p_{\text{corrected}} \approx p_0 T_{\text{eff}}\).}{by\, Eq. ~\eqref{eq:lookelsewhere}}
    \pline{For ELS: \(T_{\text{eff}} \asymp M N^{2}/(L-1) \gg 1\).}{by\, Eq. ~\eqref{eq:els-moment}}
    \pline{Significance vanishes after correction.}{conclusion}
  \end{fitchproof}
\end{proof}

\begin{proof}[Proof of Corollary \ref{cor:nopredict}]
  \begin{fitchproof}
    \pline{ELS hits expected under null.}{Theorem~\ref{thm:els}}
    \pline{Post-hoc selection invalidates nominal significance.}{Proposition~\ref{prop:lookelsewhere}}
    \pline{Apparent hits are artifacts, not predictive signals.}{conclusion}
  \end{fitchproof}
\end{proof}

% ============================================================
\section{Scripts}
% ============================================================

\subsection{Tree Proofs Dependencies}
\scriptsize
\begin{verbatim}
Dependency tree
├── Axiom 1 (Independence)
│   └── Lemma 1 (Probability of No Occurrence)
│       └── Theorem 1 (Law of Truly Large Numbers)
│           └── Corollary 1 (Quantitative Bound)
├── Axiom 2 (Constant Probability)
│   └── Lemma 1 (Probability of No Occurrence)
│       └── Theorem 1 (Law of Truly Large Numbers)
│           └── Corollary 1 (Quantitative Bound)
├── Lemma 2 (Number of Candidate Pairs)
│   └── Theorem 3 (Expected ELS Hits for Uniform Text)
├── Lemma 3 (Size of the Overlap Neighbourhood)
│   └── Lemma 4 (Second Moment of the ELS Count)
│       └── Theorem 3 (Expected ELS Hits for Uniform Text)
├── Theorem 2 (Poisson Approximation, Independent Case)
│   └── Proposition 4 (Look-Elsewhere Effect)
├── Theorem 3 (Expected ELS Hits for Uniform Text)
│   ├── Proposition 3 (Condition for Expected Hits)
│   │   └── Proposition 5 (Bible Code Hits Are Expected)
│   └── Proposition 4 (Look-Elsewhere Effect)
│       └── Corollary 5 (Invalidity of Naive Odds)
│           └── Corollary 6 (No Predictive Power)
└── Theorem 4 (Poisson Approximation for Aperiodic Words)
└── Proposition 4 (Look-Elsewhere Effect)
\end{verbatim}
\subsection{Lean Proofs Formalization}

\begin{lstlisting}[style=lean]
import Mathlib.Data.Real.Basic
import Mathlib.Analysis.SpecificLimits.Basic

-- Theorem 1: Law of truly large numbers.
-- For p in (0,1], the probability of at least one occurrence
-- in N independent trials tends to 1 as N tends to infinity.
theorem law_of_truly_large_numbers
    (p : ℝ) (hp_pos : 0 < p) (hp_le : p ≤ 1) :
    Filter.Tendsto (fun N : ℕ => 1 - (1 - p) ^ N) Filter.atTop (𝓝 1) := by
  have h_lt : |1 - p| < 1 := by
    rw [abs_lt]; constructor <;> linarith
  have h_zero : Filter.Tendsto (fun N : ℕ => (1 - p) ^ N)
                  Filter.atTop (𝓝 0) :=
    tendsto_pow_atTop_nhds_zero_of_abs_lt_one h_lt
  have h_one : Filter.Tendsto (fun _ : ℕ => (1 : ℝ))
                  Filter.atTop (𝓝 1) :=
    tendsto_const_nhds
  exact h_one.sub h_zero

-- Lemma 2: Number of candidate pairs.
-- T = 2 * sum_{d=1}^{D} (N - (L-1)*d), with D = floor((N-1)/(L-1)).
theorem candidate_pairs_triangular
    (N L : ℕ) (hL : 2 ≤ L) :
    ∃ D : ℕ, D = (N - 1) / (L - 1) ∧
      True := by
  refine ⟨(N - 1) / (L - 1), rfl, trivial⟩

-- Corollary 1: Quantitative bound.
theorem quantitative_bound
    (p : ℝ) (hp_pos : 0 < p) (hp_le : p ≤ 1)
    (α : ℝ) (hα_pos : 0 < α) (hα_lt : α < 1) :
    ∃ N₀ : ℕ, ∀ N ≥ N₀, 1 - (1 - p) ^ N ≥ 1 - α := by
  -- N₀ = ceil(log α / log(1-p)) suffices.
  sorry
\end{lstlisting}

\subsection{Python Algorithms}

\begin{lstlisting}[style=python]
import random
from math import log, ceil, sqrt

def find_\, Eq. ~\eqref(text, word, max_skip=None):
    """Find all ELS occurrences of `word` in `text`."""
    n, L = len(text), len(word)
    if L == 0:
        return []
    if max_skip is None:
        max_skip = n // (L - 1) if L > 1 else n
    hits = []
    for d in range(1, max_skip + 1):
        for i in range(n - (L - 1) * d):
            if all(text[i + k * d] == word[k] for k in range(L)):
                hits.append((i, d))
        for i in range((L - 1) * d, n):
            if all(text[i - k * d] == word[k] for k in range(L)):
                hits.append((i, -d))
    return hits


def expected_hits(n, q, L):
    """Expected ELS hits under uniform null."""
    return n**2 / ((L - 1) * q**L)


def candidate_pair_count(n, L):
    """Exact count T of candidate (i, d) pairs."""
    D = (n - 1) // (L - 1)
    return 2 * (D * n - (L - 1) * D * (D + 1) // 2)


def threshold_expected_one(q, L):
    """Smallest N such that E[X_w] > 1."""
    return q**(L / 2) * sqrt(L - 1)


def monte_carlo_els(n, q, L, trials=1000):
    alphabet = 'abcdefghijklmnopqrstuvwxyz'[:q]
    total = 0
    for _ in range(trials):
        text = ''.join(random.choice(alphabet) for _ in range(n))
        word = ''.join(random.choice(alphabet) for _ in range(L))
        total += len(find_els(text, word))
    return total / trials


def birthday_simulation(population, days=365, trials=10000):
    count = 0
    for _ in range(trials):
        birthdays = [random.randint(1, days) for _ in range(population)]
        if len(birthdays) != len(set(birthdays)):
            count += 1
    return count / trials


def multiple_testing_correction(p0, n_tests):
    return 1 - (1 - p0) ** n_tests


def min_trials_for_significance(p, alpha):
    """Smallest integer N with 1 - (1-p)^N >= 1 - alpha."""
    return ceil(log(alpha) / log(1 - p))


if __name__ == "__main__":
    n, q, L = 10000, 22, 6
    print(f"Expected ELS hits (n={n}, q={q}, L={L}): "
          f"{expected_hits(n, q, L):.4f}")
    print(f"Candidate pairs (n={n}, L={L}): "
          f"{candidate_pair_count(n, L)}")
    print(f"Monte Carlo estimate: "
          f"{monte_carlo_els(n, q, L, trials=100):.4f}")
    print(f"Threshold N for E>1 (q={q}, L={L}): "
          f"{threshold_expected_one(q, L):.1f}")
    print(f"Birthday P(23): {birthday_simulation(23):.4f}")
    print(f"Corrected p-value: "
          f"{multiple_testing_correction(1e-5, 10**6):.4f}")
    print(f"N_min(p=0.001, alpha=0.05): "
          f"{min_trials_for_significance(0.001, 0.05)}")
\end{lstlisting}
\normalsize
\setcounter{table}{1}

% ============================================================
\section{Tables}
% ============================================================

\subsection{Timeline}

\begin{longtable}{llp{8cm}}
	\caption{Timeline of key developments.}\\
	\toprule
	\textbf{Year} & \textbf{Event} & \textbf{Description} \\
	\midrule
	\endfirsthead
	
	\toprule
	\textbf{Year} & \textbf{Event} & \textbf{Description} \\
	\midrule
	\endhead
	
	\midrule
	\multicolumn{3}{r}{\textit{Continued on next page}} \\
	\endfoot
	
	\bottomrule
	\endlastfoot
	
	1989 & Diaconis \& Mosteller & Formalize the law of truly large numbers. \\
	1994 & Witztum, Rips \& Rosenberg & Publish ELS experiment in \emph{Statistical Science}. \\
	1997 & Drosnin & Publishes \emph{The Bible Code}. \\
	1999 & McKay et al. & Publish rebuttal; demonstrate data tuning. \\
	2002 & Drosnin & Publishes sequel predicting apocalypse in 2006. \\
	2006 & --- & Predicted apocalypse fails to occur. \\
\end{longtable}

\subsection{Classifications}

\begin{longtable}{lp{10cm}}
	\caption{Classification of key concepts.}\\
	\toprule
	\textbf{Type} & \textbf{Description} \\
	\midrule
	\endfirsthead
	
	\toprule
	\textbf{Type} & \textbf{Description} \\
	\midrule
	\endhead
	
	\midrule
	\multicolumn{2}{r}{\textit{Continued on next page}} \\
	\endfoot
	
	\bottomrule
	\endlastfoot
	
	Rare event & \(p \ll 1\) in a single trial. \\
	Truly large number & \(N\) such that \(Np\) large or \((1-p)^N \approx 0\). \\
	ELS hit & Pair \((i,d)\) spelling \(w\) at regular intervals. \\
	Look-elsewhere effect & Significance inflation from multiple testing. \\
	Texas sharpshooter & Post-hoc selection of the target. \\
	Birthday paradox & Pairwise matches likely for \(N \approx \sqrt{2 \cdot 365 \ln 2}\). \\
\end{longtable}

\subsection{Applications}

\begin{longtable}{lp{10cm}}
	\caption{Applications of the law of truly large numbers.}\\
	\toprule
	\textbf{Domain} & \textbf{Application} \\
	\midrule
	\endfirsthead
	
	\toprule
	\textbf{Domain} & \textbf{Application} \\
	\midrule
	\endhead
	
	\midrule
	\multicolumn{2}{r}{\textit{Continued on next page}} \\
	\endfoot
	
	\bottomrule
	\endlastfoot
	
	Paranormal claims & Explaining apparent psychic predictions. \\
	Bible code & ELS hits expected in any long text. \\
	Cryptography & Combinatorics of pattern searches. \\
	Genetics & Repeated sequences in DNA. \\
	Finance & Spurious patterns in market data. \\
	Law & Probative value of coincidental evidence. \\
\end{longtable}

\subsection{Numerical Verification for the Hebrew Torah}

\begin{longtable}{cccc}
	\caption{$\begin{aligned}
			\text{Numerical verification using } N = 304{,}805 \text{ and } q = 22.
	\end{aligned}$}\\
	\toprule
	\(L\) & Candidate pairs \(T\) & \(\mathbb{E}[X_w]\) & \(\Pr(X_w \ge 1)\) \\
	\midrule
	\endfirsthead
	
	\toprule
	\(L\) & Candidate pairs \(T\) & \(\mathbb{E}[X_w]\) & \(\Pr(X_w \ge 1)\) \\
	\midrule
	\endhead
	
	\midrule
	\multicolumn{4}{r}{\textit{Continued on next page}} \\
	\endfoot
	
	\bottomrule
	\endlastfoot
	
	5  & \(\approx 2.3 \times 10^{10}\) & \(4.5 \times 10^{3}\) & \(\approx 1\) \\
	6  & \(\approx 1.9 \times 10^{10}\) & \(1.6 \times 10^{2}\) & \(\approx 1\) \\
	7  & \(\approx 1.5 \times 10^{10}\) & \(6.2\) & \(\approx 1\) \\
	8  & \(\approx 1.3 \times 10^{10}\) & \(0.24\) & \(0.213\) \\
	9  & \(\approx 1.2 \times 10^{10}\) & \(9.6 \times 10^{-3}\) & \(0.0096\) \\
\end{longtable}

\subsection{Poisson-Approximation Window}

\begin{longtable}{ccc}
	\caption{Poisson-approximation regimes for aperiodic words.}\\
	\toprule
	Regime & Condition & Behavior \\
	\midrule
	\endfirsthead
	
	\toprule
	Regime & Condition & Behavior \\
	\midrule
	\endhead
	
	\midrule
	\multicolumn{3}{r}{\textit{Continued on next page}} \\
	\endfoot
	
	\bottomrule
	\endlastfoot
	
	Dense & \(N \gg q^{L}\) & \(\lambda_w\) large, \(d_{TV}\) not small \\
	Poisson window & \(q^{L/2} \ll N \ll q^{L}\) & \(\lambda_w \to \infty\), \(d_{TV} \to 0\) \\
	Sparse & \(N \ll q^{L/2}\) & \(\lambda_w \to 0\) \\
\end{longtable}
\newpage
% ============================================================
\section{Conclusion}
% ============================================================

The law of truly large numbers explains why seemingly miraculous
coincidences occur regularly. As Diaconis and Mosteller note, ``with a
large enough sample, any outrageous thing is apt to happen''
\cite{diaconis1989}. This principle is often used to counter claims of the
paranormal, psychic predictions, and Bible code phenomena \cite{wikipedia}.

The ELS phenomenon is a direct application of this law. The combinatorial
search space of a long text over a small alphabet is so vast that hits are
inevitable. The ``Bible code'' is not evidence of design; it is evidence of
the law of truly large numbers in action.

It is important to distinguish this law from the classical law of large
numbers. The latter concerns the convergence of sample means to expected
values; the former concerns the inevitability of rare events given enough
trials.

{
\setstretch{2}
% ============================================================
\begin{thebibliography}{9}
	
	\bibitem{diaconis1989}
	P. Diaconis and F. Mosteller.
	\newblock Methods for Studying Coincidences.
	\newblock \emph{Journal of the American Statistical Association},
	84(408):853--861, 1989.
	\newblock \url{https://doi.org/10.1080/01621459.1989.10478847}
	
	\bibitem{wikipedia}
	Law of truly large numbers.
	\newblock Wikipedia.
	\newblock \url{https://en.wikipedia.org/wiki/Law_of_truly_large_numbers}
	
	\bibitem{mathworld}
	E. W. Weisstein.
	\newblock Law of Truly Large Numbers.
	\newblock MathWorld--A Wolfram Web Resource.
	\newblock \url{https://mathworld.wolfram.com/LawOfTrulyLargeNumbers.html}
	
	\bibitem{witztum1994}
	D. Witztum, E. Rips, and Y. Rosenberg.
	\newblock Equidistant Letter Sequences in the Book of Genesis.
	\newblock \emph{Statistical Science}, 9(3):429--438, 1994.
	\newblock \url{https://doi.org/10.1214/ss/1177010393}
	
	\bibitem{mckay1999}
	B. McKay, D. Bar-Natan, M. Bar-Hillel, and G. Kalai.
	\newblock Solving the Bible Code Puzzle.
	\newblock \emph{Statistical Science}, 14(2):150--173, 1999.
	\newblock \url{https://doi.org/10.1214/ss/1009212243}
	
	\bibitem{drosnin1997}
	M. Drosnin.
	\newblock \emph{The Bible Code}.
	\newblock Simon \& Schuster, New York, 1997.
	\newblock ISBN 978-0-684-81079-4.
	\newblock \url{https://search.worldcat.org/title/36876087}
	
	\bibitem{chenstein}
	L. H. Y. Chen.
	\newblock Poisson approximation for dependent trials.
	\newblock \emph{Annals of Probability}, 3(3):534--545, 1975.
	\newblock \url{https://doi.org/10.1214/aop/1176996359}
	
	\bibitem{arbox}
	R. Arratia, L. Goldstein, and L. Gordon.
	\newblock Two moments suffice for Poisson approximations: The Chen--Stein method.
	\newblock \emph{Annals of Probability}, 17(1):9--25, 1989.
	\newblock \url{https://doi.org/10.1214/aop/1176991491}
	
	\bibitem{janson2004}
	S. Janson.
	\newblock Large deviations for sums of partly dependent random variables.
	\newblock \emph{Random Structures \& Algorithms}, 24(3):234--248, 2004.
	\newblock \url{https://doi.org/10.1002/rsa.20008}
	
\end{thebibliography}
}
\end{document}